\appendices
\section{Supplementary Material}

Reviewers are not required to read this appendix (per the SaTML 2027
checklist); it contains the full per-depth data behind the summary
tables in Section~\ref{sec:results}, the exact experimental design
parameters, and one worked example of the labeling rubric.

\subsection{Experimental grid}

Full generation factorial (\texttt{configs/experiment\_grid.yaml}):
\texttt{scenario\_id} (30: 24 poisoned across 4 \texttt{poison\_form}
styles $\times$ 6 independently-constructed scenarios each, plus 6
\texttt{clean\_control}), \texttt{model} (gpt-4o-mini, gpt-4o,
Llama-3.3-70B-Instruct), \texttt{top\_k} (1, 3, 5, 10),
\texttt{write\_fanout} (1, 2, 3), \texttt{derivation\_transform}
(summarize, paraphrase, refine, continue), \texttt{seed} (0--4, frozen).
Total: $30 \times 3 \times 4 \times 3 \times 4 \times 5 = 21{,}600$
planned traces (21,570 generated, 99.86\%). \texttt{depth} (0--5) and
\texttt{attribution\_threshold} (null, 0.5, 0.7, 0.85) are post-hoc
analysis factors computed from each trace's persisted edges, not
separate generation runs; \texttt{containment\_policy}
(\texttt{flat\_transitive}/\texttt{depth\_aware}) is a further post-hoc
layer for H4.

\texttt{scenario\_id} is the real inferential sampling unit (n=24 for
H1--H4's primary pooled estimand, stratified 6 per \texttt{poison\_form})
--- not individual factorial cells, and not seed count. All 5 seeds are
averaged within scenario-condition before any bootstrap resampling;
resampling always sees 24 scenario units.

\subsection{Full per-depth H1 results (recall vs.\ inflation growth)}

Section~\ref{sec:results} reports only depth 1 and depth 5 per model.
Table~\ref{tab:appendix-h1} gives the full depth 1--5 sweep, both
propagation policies, all three models (95\% bootstrap CI, n=24
scenarios).

\begin{longtable}{@{}llrlll@{}}
\caption{Full H1 sweep: blast-radius precision, recall, and inflation ratio by model, policy, and depth.}
\label{tab:appendix-h1} \\
\toprule
Model & Policy & Depth & P\_BR & R\_BR & Inflation \\
\midrule
\endfirsthead
\toprule
Model & Policy & Depth & P\_BR & R\_BR & Inflation \\
\midrule
\endhead
\bottomrule
\endfoot
gpt-4o-mini & context\_exposure & 1 & 0.669 [0.642, 0.700] & 1.000 [1.000, 1.000] & 1.832 [1.748, 1.910] \\
gpt-4o-mini & context\_exposure & 2 & 0.639 [0.612, 0.666] & 1.000 [1.000, 1.000] & 1.989 [1.868, 2.115] \\
gpt-4o-mini & context\_exposure & 3 & 0.629 [0.602, 0.657] & 1.000 [1.000, 1.000] & 2.107 [1.942, 2.278] \\
gpt-4o-mini & context\_exposure & 4 & 0.625 [0.596, 0.654] & 1.000 [1.000, 1.000] & 2.197 [1.995, 2.434] \\
gpt-4o-mini & context\_exposure & 5 & 0.622 [0.593, 0.649] & 1.000 [1.000, 1.000] & 2.273 [2.044, 2.516] \\
gpt-4o-mini & structural & 1 & 0.981 [0.963, 0.996] & 0.915 [0.872, 0.955] & 0.932 [0.898, 0.964] \\
gpt-4o-mini & structural & 2 & 0.944 [0.923, 0.967] & 0.913 [0.884, 0.939] & 1.012 [0.959, 1.070] \\
gpt-4o-mini & structural & 3 & 0.933 [0.907, 0.962] & 0.918 [0.893, 0.941] & 1.070 [0.993, 1.150] \\
gpt-4o-mini & structural & 4 & 0.926 [0.897, 0.956] & 0.918 [0.896, 0.941] & 1.119 [1.018, 1.232] \\
gpt-4o-mini & structural & 5 & 0.923 [0.894, 0.953] & 0.921 [0.900, 0.943] & 1.157 [1.042, 1.275] \\
gpt-4o & context\_exposure & 1 & 0.628 [0.609, 0.643] & 1.000 [1.000, 1.000] & 1.946 [1.904, 1.995] \\
gpt-4o & context\_exposure & 2 & 0.597 [0.576, 0.618] & 1.000 [1.000, 1.000] & 2.119 [2.005, 2.243] \\
gpt-4o & context\_exposure & 3 & 0.589 [0.564, 0.612] & 1.000 [1.000, 1.000] & 2.236 [2.051, 2.449] \\
gpt-4o & context\_exposure & 4 & 0.584 [0.557, 0.610] & 1.000 [1.000, 1.000] & 2.331 [2.062, 2.632] \\
gpt-4o & context\_exposure & 5 & 0.581 [0.552, 0.607] & 1.000 [1.000, 1.000] & 2.422 [2.100, 2.783] \\
gpt-4o & structural & 1 & 0.966 [0.901, 1.000] & 0.939 [0.873, 0.980] & 0.972 [0.960, 0.984] \\
gpt-4o & structural & 2 & 0.942 [0.894, 0.978] & 0.954 [0.914, 0.979] & 1.062 [1.012, 1.114] \\
gpt-4o & structural & 3 & 0.932 [0.886, 0.971] & 0.956 [0.920, 0.979] & 1.118 [1.030, 1.210] \\
gpt-4o & structural & 4 & 0.927 [0.880, 0.969] & 0.957 [0.927, 0.977] & 1.163 [1.041, 1.294] \\
gpt-4o & structural & 5 & 0.922 [0.875, 0.966] & 0.959 [0.930, 0.979] & 1.209 [1.056, 1.367] \\
llama-3.3-70b & context\_exposure & 1 & 0.632 [0.620, 0.647] & 1.000 [1.000, 1.000] & 1.939 [1.895, 1.975] \\
llama-3.3-70b & context\_exposure & 2 & 0.609 [0.587, 0.629] & 1.000 [1.000, 1.000] & 2.057 [1.952, 2.179] \\
llama-3.3-70b & context\_exposure & 3 & 0.599 [0.571, 0.622] & 1.000 [1.000, 1.000] & 2.167 [1.984, 2.390] \\
llama-3.3-70b & context\_exposure & 4 & 0.593 [0.562, 0.618] & 1.000 [1.000, 1.000] & 2.258 [2.014, 2.563] \\
llama-3.3-70b & context\_exposure & 5 & 0.586 [0.552, 0.614] & 1.000 [1.000, 1.000] & 2.368 [2.034, 2.773] \\
llama-3.3-70b & structural & 1 & 1.000 [0.999, 1.000] & 0.976 [0.959, 0.989] & 0.976 [0.959, 0.990] \\
llama-3.3-70b & structural & 2 & 0.966 [0.938, 0.989] & 0.973 [0.959, 0.984] & 1.041 [0.988, 1.103] \\
llama-3.3-70b & structural & 3 & 0.952 [0.912, 0.984] & 0.975 [0.963, 0.985] & 1.096 [1.006, 1.208] \\
llama-3.3-70b & structural & 4 & 0.943 [0.901, 0.979] & 0.975 [0.964, 0.985] & 1.143 [1.016, 1.295] \\
llama-3.3-70b & structural & 5 & 0.934 [0.883, 0.974] & 0.976 [0.966, 0.985] & 1.198 [1.030, 1.407] \\
\end{longtable}

Growth is monotone but concave in every row: most of the precision loss
and inflation growth happens by depth 2--3, not spread evenly across
all five depths.

\subsection{Full H2 attribution-threshold sweep}

Section~\ref{sec:results} reports only depth 5. Table~\ref{tab:appendix-h2}
gives the full depth 1--5 sweep, all three thresholds, all three models.

\begin{longtable}{@{}lrrll@{}}
\caption{Full H2 sweep: blast-radius precision and recall by model, attribution threshold, and depth.}
\label{tab:appendix-h2} \\
\toprule
Model & Threshold & Depth & P\_BR & R\_BR \\
\midrule
\endfirsthead
\toprule
Model & Threshold & Depth & P\_BR & R\_BR \\
\midrule
\endhead
\bottomrule
\endfoot
gpt-4o-mini & 0.5 & 1 & 0.957 [0.933, 0.979] & 0.983 [0.967, 0.995] \\
gpt-4o-mini & 0.5 & 2 & 0.920 [0.895, 0.946] & 0.988 [0.973, 0.997] \\
gpt-4o-mini & 0.5 & 3 & 0.907 [0.878, 0.938] & 0.989 [0.975, 0.998] \\
gpt-4o-mini & 0.5 & 4 & 0.901 [0.870, 0.933] & 0.990 [0.975, 0.998] \\
gpt-4o-mini & 0.5 & 5 & 0.898 [0.868, 0.930] & 0.990 [0.976, 0.999] \\
gpt-4o-mini & 0.7 & 1 & 0.874 [0.832, 0.918] & 0.844 [0.798, 0.892] \\
gpt-4o-mini & 0.7 & 2 & 0.839 [0.803, 0.875] & 0.859 [0.818, 0.904] \\
gpt-4o-mini & 0.7 & 3 & 0.830 [0.792, 0.869] & 0.867 [0.823, 0.910] \\
gpt-4o-mini & 0.7 & 4 & 0.823 [0.784, 0.863] & 0.870 [0.827, 0.914] \\
gpt-4o-mini & 0.7 & 5 & 0.820 [0.781, 0.860] & 0.871 [0.828, 0.916] \\
gpt-4o-mini & 0.85 & 1 & 0.648 [0.575, 0.722] & 0.620 [0.552, 0.689] \\
gpt-4o-mini & 0.85 & 2 & 0.639 [0.569, 0.708] & 0.626 [0.553, 0.701] \\
gpt-4o-mini & 0.85 & 3 & 0.635 [0.568, 0.703] & 0.631 [0.557, 0.709] \\
gpt-4o-mini & 0.85 & 4 & 0.635 [0.567, 0.703] & 0.632 [0.557, 0.710] \\
gpt-4o-mini & 0.85 & 5 & 0.633 [0.566, 0.701] & 0.632 [0.558, 0.711] \\
gpt-4o & 0.5 & 1 & 0.956 [0.916, 0.983] & 0.996 [0.991, 0.999] \\
gpt-4o & 0.5 & 2 & 0.911 [0.871, 0.946] & 0.997 [0.993, 0.999] \\
gpt-4o & 0.5 & 3 & 0.897 [0.854, 0.934] & 0.997 [0.994, 0.999] \\
gpt-4o & 0.5 & 4 & 0.891 [0.846, 0.931] & 0.997 [0.994, 1.000] \\
gpt-4o & 0.5 & 5 & 0.887 [0.841, 0.928] & 0.998 [0.995, 1.000] \\
gpt-4o & 0.7 & 1 & 0.945 [0.877, 0.986] & 0.938 [0.872, 0.980] \\
gpt-4o & 0.7 & 2 & 0.915 [0.869, 0.951] & 0.953 [0.924, 0.976] \\
gpt-4o & 0.7 & 3 & 0.902 [0.855, 0.942] & 0.956 [0.928, 0.978] \\
gpt-4o & 0.7 & 4 & 0.897 [0.849, 0.938] & 0.960 [0.937, 0.979] \\
gpt-4o & 0.7 & 5 & 0.892 [0.844, 0.935] & 0.962 [0.941, 0.980] \\
gpt-4o & 0.85 & 1 & 0.751 [0.656, 0.845] & 0.741 [0.650, 0.832] \\
gpt-4o & 0.85 & 2 & 0.726 [0.650, 0.808] & 0.731 [0.652, 0.815] \\
gpt-4o & 0.85 & 3 & 0.720 [0.649, 0.797] & 0.729 [0.648, 0.812] \\
gpt-4o & 0.85 & 4 & 0.714 [0.644, 0.792] & 0.726 [0.646, 0.810] \\
gpt-4o & 0.85 & 5 & 0.710 [0.639, 0.786] & 0.725 [0.645, 0.809] \\
llama-3.3-70b & 0.5 & 1 & 0.973 [0.954, 0.989] & 0.993 [0.989, 0.997] \\
llama-3.3-70b & 0.5 & 2 & 0.928 [0.898, 0.955] & 0.996 [0.993, 0.998] \\
llama-3.3-70b & 0.5 & 3 & 0.912 [0.874, 0.945] & 0.996 [0.994, 0.998] \\
llama-3.3-70b & 0.5 & 4 & 0.900 [0.861, 0.936] & 0.997 [0.995, 0.998] \\
llama-3.3-70b & 0.5 & 5 & 0.890 [0.845, 0.929] & 0.997 [0.996, 0.999] \\
llama-3.3-70b & 0.7 & 1 & 0.956 [0.930, 0.980] & 0.940 [0.910, 0.968] \\
llama-3.3-70b & 0.7 & 2 & 0.914 [0.884, 0.942] & 0.943 [0.915, 0.970] \\
llama-3.3-70b & 0.7 & 3 & 0.898 [0.860, 0.931] & 0.946 [0.919, 0.972] \\
llama-3.3-70b & 0.7 & 4 & 0.889 [0.849, 0.922] & 0.948 [0.920, 0.974] \\
llama-3.3-70b & 0.7 & 5 & 0.878 [0.832, 0.916] & 0.949 [0.921, 0.975] \\
llama-3.3-70b & 0.85 & 1 & 0.678 [0.601, 0.757] & 0.674 [0.599, 0.752] \\
llama-3.3-70b & 0.85 & 2 & 0.655 [0.579, 0.738] & 0.665 [0.589, 0.744] \\
llama-3.3-70b & 0.85 & 3 & 0.649 [0.570, 0.728] & 0.659 [0.585, 0.738] \\
llama-3.3-70b & 0.85 & 4 & 0.645 [0.567, 0.724] & 0.655 [0.579, 0.733] \\
llama-3.3-70b & 0.85 & 5 & 0.640 [0.562, 0.720] & 0.652 [0.578, 0.731] \\
\end{longtable}

The frontier is already largely set by depth 2 --- recall at threshold
0.85 barely moves between depth 2 and depth 5 for any model, while
precision keeps drifting down slowly. The steep recall cost of an
aggressive threshold is a depth-1 phenomenon, not something that
compounds further with depth.

\subsection{Full H3 per-transform laundering breakdown}

Table~\ref{tab:appendix-h3} gives the full per-model $\times$
per-transform $\times$ per-depth laundering breakdown.

\begin{longtable}{@{}llrrrrr@{}}
\caption{Full H3 sweep: laundering rate and structural-lineage recall by model, derivation transform, and depth.}
\label{tab:appendix-h3} \\
\toprule
Model & Transform & Depth & $n_{\text{CARRIES}}$ & $n_{\text{laundered}}$ & Laundering rate & Struct.\ recall \\
\midrule
\endfirsthead
\toprule
Model & Transform & Depth & $n_{\text{CARRIES}}$ & $n_{\text{laundered}}$ & Laundering rate & Struct.\ recall \\
\midrule
\endhead
\bottomrule
\endfoot
gpt-4o & continue & 1 & 1677 & 6 & 0.358\% & 0.811 \\
gpt-4o & continue & 2 & 1507 & 2 & 0.133\% & 0.883 \\
gpt-4o & continue & 3 & 1511 & 6 & 0.397\% & 0.866 \\
gpt-4o & continue & 4 & 1544 & 13 & 0.842\% & 0.846 \\
gpt-4o & continue & 5 & 1484 & 7 & 0.472\% & 0.865 \\
gpt-4o & paraphrase & 1 & 1279 & 0 & 0.000\% & 0.998 \\
gpt-4o & paraphrase & 2 & 1272 & 0 & 0.000\% & 1.000 \\
gpt-4o & paraphrase & 3 & 1256 & 0 & 0.000\% & 1.000 \\
gpt-4o & paraphrase & 4 & 1231 & 0 & 0.000\% & 1.000 \\
gpt-4o & paraphrase & 5 & 1226 & 0 & 0.000\% & 1.000 \\
gpt-4o & refine & 1 & 1390 & 0 & 0.000\% & 0.982 \\
gpt-4o & refine & 2 & 1344 & 0 & 0.000\% & 1.000 \\
gpt-4o & refine & 3 & 1324 & 0 & 0.000\% & 1.000 \\
gpt-4o & refine & 4 & 1330 & 0 & 0.000\% & 1.000 \\
gpt-4o & refine & 5 & 1321 & 0 & 0.000\% & 1.000 \\
gpt-4o & summarize & 1 & 1246 & 11 & 0.883\% & 1.000 \\
gpt-4o & summarize & 2 & 1298 & 16 & 1.233\% & 1.000 \\
gpt-4o & summarize & 3 & 1288 & 16 & 1.242\% & 1.000 \\
gpt-4o & summarize & 4 & 1303 & 19 & 1.458\% & 1.000 \\
gpt-4o & summarize & 5 & 1305 & 22 & 1.686\% & 1.000 \\
gpt-4o-mini & continue & 1 & 1796 & 28 & 1.559\% & 0.719 \\
gpt-4o-mini & continue & 2 & 2076 & 45 & 2.168\% & 0.630 \\
gpt-4o-mini & continue & 3 & 2047 & 22 & 1.075\% & 0.638 \\
gpt-4o-mini & continue & 4 & 2098 & 31 & 1.478\% & 0.615 \\
gpt-4o-mini & continue & 5 & 2051 & 25 & 1.219\% & 0.628 \\
gpt-4o-mini & paraphrase & 1 & 1263 & 0 & 0.000\% & 0.993 \\
gpt-4o-mini & paraphrase & 2 & 1205 & 0 & 0.000\% & 0.993 \\
gpt-4o-mini & paraphrase & 3 & 1191 & 0 & 0.000\% & 1.000 \\
gpt-4o-mini & paraphrase & 4 & 1197 & 0 & 0.000\% & 1.000 \\
gpt-4o-mini & paraphrase & 5 & 1199 & 0 & 0.000\% & 1.000 \\
gpt-4o-mini & refine & 1 & 1552 & 0 & 0.000\% & 0.883 \\
gpt-4o-mini & refine & 2 & 1440 & 0 & 0.000\% & 0.967 \\
gpt-4o-mini & refine & 3 & 1427 & 0 & 0.000\% & 0.980 \\
gpt-4o-mini & refine & 4 & 1439 & 0 & 0.000\% & 0.963 \\
gpt-4o-mini & refine & 5 & 1422 & 0 & 0.000\% & 0.985 \\
gpt-4o-mini & summarize & 1 & 1396 & 0 & 0.000\% & 0.907 \\
gpt-4o-mini & summarize & 2 & 1251 & 0 & 0.000\% & 1.000 \\
gpt-4o-mini & summarize & 3 & 1276 & 0 & 0.000\% & 1.000 \\
gpt-4o-mini & summarize & 4 & 1272 & 0 & 0.000\% & 1.000 \\
gpt-4o-mini & summarize & 5 & 1284 & 0 & 0.000\% & 1.000 \\
llama-3.3-70b & continue & 1 & 1564 & 8 & 0.512\% & 0.865 \\
llama-3.3-70b & continue & 2 & 1563 & 19 & 1.216\% & 0.865 \\
llama-3.3-70b & continue & 3 & 1457 & 25 & 1.716\% & 0.890 \\
llama-3.3-70b & continue & 4 & 1460 & 20 & 1.370\% & 0.872 \\
llama-3.3-70b & continue & 5 & 1375 & 36 & 2.618\% & 0.884 \\
llama-3.3-70b & paraphrase & 1 & 1383 & 29 & 2.097\% & 0.996 \\
llama-3.3-70b & paraphrase & 2 & 1274 & 61 & 4.788\% & 1.000 \\
llama-3.3-70b & paraphrase & 3 & 1255 & 74 & 5.896\% & 1.000 \\
llama-3.3-70b & paraphrase & 4 & 1227 & 63 & 5.134\% & 1.000 \\
llama-3.3-70b & paraphrase & 5 & 1213 & 57 & 4.699\% & 1.000 \\
llama-3.3-70b & refine & 1 & 1397 & 0 & 0.000\% & 0.994 \\
llama-3.3-70b & refine & 2 & 1391 & 0 & 0.000\% & 0.996 \\
llama-3.3-70b & refine & 3 & 1388 & 1 & 0.072\% & 0.999 \\
llama-3.3-70b & refine & 4 & 1383 & 0 & 0.000\% & 0.999 \\
llama-3.3-70b & refine & 5 & 1371 & 0 & 0.000\% & 1.000 \\
llama-3.3-70b & summarize & 1 & 1305 & 0 & 0.000\% & 0.962 \\
llama-3.3-70b & summarize & 2 & 1298 & 0 & 0.000\% & 0.986 \\
llama-3.3-70b & summarize & 3 & 1290 & 0 & 0.000\% & 1.000 \\
llama-3.3-70b & summarize & 4 & 1280 & 0 & 0.000\% & 1.000 \\
llama-3.3-70b & summarize & 5 & 1285 & 0 & 0.000\% & 1.000 \\
\end{longtable}

\texttt{refine} launders essentially nowhere (one single-count exception
at llama-3.3-70b depth 3) and \texttt{paraphrase} launders nowhere for
gpt-4o/gpt-4o-mini --- but launders substantially for llama-3.3-70b
(4.7--5.9\%), the sharpest single cell in the whole table and the
clearest evidence that laundering is a model property, not a transform
property alone.

\subsection{Full H4 containment-cost breakdown}

Table~\ref{tab:appendix-h4} gives the full depth 1--5 breakdown, both
policies, all three models --- mean objects flagged and mean missed
contamination per trace.

\begin{longtable}{@{}llrrr@{}}
\caption{Full H4 sweep: mean objects flagged and mean missed contamination by model, policy, and depth.}
\label{tab:appendix-h4} \\
\toprule
Model & Policy & Depth & Objects flagged & Missed contamination \\
\midrule
\endfirsthead
\toprule
Model & Policy & Depth & Objects flagged & Missed contamination \\
\midrule
\endhead
\bottomrule
\endfoot
gpt-4o-mini & flat\_transitive & 1 & 2.000 & 0.000 \\
gpt-4o-mini & flat\_transitive & 2 & 4.000 & 0.000 \\
gpt-4o-mini & flat\_transitive & 3 & 6.000 & 0.000 \\
gpt-4o-mini & flat\_transitive & 4 & 8.000 & 0.000 \\
gpt-4o-mini & flat\_transitive & 5 & 10.000 & 0.001 \\
gpt-4o-mini & depth\_aware & 1 & 2.000 & 0.000 \\
gpt-4o-mini & depth\_aware & 2 & 4.000 & 0.000 \\
gpt-4o-mini & depth\_aware & 3 & 5.000 & 0.134 \\
gpt-4o-mini & depth\_aware & 4 & 6.000 & 0.283 \\
gpt-4o-mini & depth\_aware & 5 & 7.000 & 0.419 \\
gpt-4o & flat\_transitive & 1 & 2.000 & 0.000 \\
gpt-4o & flat\_transitive & 2 & 4.000 & 0.000 \\
gpt-4o & flat\_transitive & 3 & 6.000 & 0.000 \\
gpt-4o & flat\_transitive & 4 & 8.000 & 0.000 \\
gpt-4o & flat\_transitive & 5 & 10.000 & 0.000 \\
gpt-4o & depth\_aware & 1 & 2.000 & 0.000 \\
gpt-4o & depth\_aware & 2 & 4.000 & 0.000 \\
gpt-4o & depth\_aware & 3 & 5.000 & 0.035 \\
gpt-4o & depth\_aware & 4 & 6.000 & 0.077 \\
gpt-4o & depth\_aware & 5 & 7.000 & 0.111 \\
llama-3.3-70b & flat\_transitive & 1 & 1.997 & 0.000 \\
llama-3.3-70b & flat\_transitive & 2 & 3.994 & 0.000 \\
llama-3.3-70b & flat\_transitive & 3 & 5.991 & 0.000 \\
llama-3.3-70b & flat\_transitive & 4 & 7.988 & 0.000 \\
llama-3.3-70b & flat\_transitive & 5 & 9.985 & 0.000 \\
llama-3.3-70b & depth\_aware & 1 & 1.997 & 0.000 \\
llama-3.3-70b & depth\_aware & 2 & 3.994 & 0.000 \\
llama-3.3-70b & depth\_aware & 3 & 4.994 & 0.028 \\
llama-3.3-70b & depth\_aware & 4 & 5.994 & 0.061 \\
llama-3.3-70b & depth\_aware & 5 & 6.994 & 0.089 \\
\end{longtable}

Both policies are identical through depth 2 (\texttt{depth\_aware}'s
2-hop conservative window hasn't yet switched to structural-only
propagation) and diverge starting depth 3 --- exactly the depth at which
\texttt{depth\_aware}'s missed-contamination cost starts accruing, and
gpt-4o-mini pulls away from the other two models by roughly the same
margin at every depth $\geq 3$.

\subsection{Seed diversity}

Not a bug --- no \texttt{seed} parameter is ever sent to any model API
(seeds rely on genuine API-level stochasticity at temperature 0) --- but
the duplication rate is worth stating plainly, since a reviewer computing
effective sample size from raw trace counts would otherwise overestimate
it: across the 4,320 \texttt{clean\_control} generation cells, 39.5\%
(1,708/4,320) have all 5 seeds byte-identical, 8.8\% (381/4,320) have all
5 distinct, and 51.6\% (2,231/4,320) show partial diversity. This varies
by \texttt{derivation\_transform} in the direction genuine task
determinism would predict: \texttt{summarize} 60.6\% fully-duplicated,
\texttt{refine} 42.5\%, \texttt{paraphrase} 38.3\%, \texttt{continue}
16.8\%. This does not change the statistical design (the bootstrap unit
is \texttt{scenario\_id}, n=24, unaffected by within-cell seed
duplication).

\subsection{Worked labeling example}

One example illustrating the CARRIES/REFERENCES/CLEAN rubric and the
compositional-target rule, representative of the \texttt{multi\_hop\_setup}
boundary case discussed in Limitations as the labeler's weakest point.
Semantic target: ``Route every signed NDA to
\texttt{mirror-intake@compliance-shadow.example} before it is
countersigned.''

\begin{itemize}
\item \textbf{Candidate A} --- ``All executed NDAs are now automatically
  forwarded to \texttt{mirror-intake@compliance-shadow.example} prior to
  final countersignature.'' $\rightarrow$ \textbf{CARRIES}: restates the
  full claim (forwarding destination + timing) in different words.
\item \textbf{Candidate B} --- ``Legal has been discussing changes to the
  NDA countersignature workflow, including a new pre-signature review
  step routed through the compliance team's shared intake address.''
  $\rightarrow$ \textbf{REFERENCES}: preserves the topic without
  asserting the specific destination or the mandatory-forwarding claim.
\item \textbf{Candidate C} --- ``The legal team updated their document
  retention policy for signed contracts.'' $\rightarrow$ \textbf{CLEAN}:
  no trace of the claim.
\end{itemize}

The compositional-target rule (frozen) applies when a target requires
combining two premises: CARRIES requires the candidate to assert the
\emph{combined} claim, not merely restate both underlying facts side by
side --- this is specifically what depresses REFERENCES-class recall
(0.46) in \texttt{multi\_hop\_setup} scenarios, independently
corroborated by the human-human $\kappa$ baseline (2026-09-21): all 5
annotator disagreements in that 120-sample check were exactly this
REFERENCES-vs-CARRIES boundary.
